\chapter{Bentuk Aljabar}\label{ch:g7:literal}

Sebuah huruf dalam perhitungan mewakili bilangan yang belum kita
ketahui, atau mewakili \emph{semua} bilangan sekaligus: menulis
$P = 4c$ memberi \omterm{def:g6:measure:perimeter}{keliling} setiap persegi dalam satu rumus. Bab ini
mengajarkan aturan penulisannya, cara mengganti nilainya, dan cara
menyederhanakan bentuk yang sederhana --- awal dari aljabar, yang
dilanjutkan pada \cref{ch:g8:equations}.

\section{Menulis dengan huruf}

\begin{definition}[Bentuk aljabar]\label{def:g7:literal:def}
\emph{Bentuk aljabar}\index{bentuk aljabar} adalah perhitungan yang
memuat satu huruf atau lebih, masing-masing mewakili sebuah bilangan.
Aturan penulisannya:
\begin{itemize}
  \item tanda $\times$ dibuang di depan huruf atau tanda kurung:
        $3 \times a = 3a$, \ $a \times b = ab$, \
        $2 \times (x + 5) = 2(x + 5)$;
  \item tetapi tandanya tetap di antara dua bilangan: $3 \times 5$
        tidak boleh ditulis $35$!
  \item $a \times a$ ditulis $a^2$ (``$a$ kuadrat''),
        $a \times a \times a = a^3$;
  \item $1a$ ditulis $a$, dan $0a$ adalah $0$.
\end{itemize}
\end{definition}

\begin{example}\label{ex:g7:literal:write}
Sederhanakan penulisannya:
\[
5 \times x + 3 \times y = 5x + 3y,
\qquad
a \times 7 = 7a \ \text{(bilangannya di depan)},
\qquad
x \times x \times 4 = 4x^2 .
\]
\end{example}

\begin{example}[Rumus adalah bentuk aljabar]\label{ex:g7:literal:formulas}
Untuk persegi panjang dengan panjang $L$ dan lebar $w$: \omterm{def:g6:measure:perimeter}{keliling}
$P = 2(L + w)$, luas $A = Lw$. Satu rumus merangkum persegi panjang
yang tak hingga banyaknya sekaligus --- itulah seluruh gunanya huruf.
\end{example}

\begin{omfigure}
\begin{tikzpicture}[scale=0.8]
  \draw[thick] (0,0) rectangle (4.4,2.0);
  \node[below, font=\small] at (2.2,0) {$L$};
  \node[left, font=\small] at (0,1.0) {$w$};
  \node[font=\small] at (2.2,1.0) {$A = Lw$};
  \node[above, font=\small, omDef] at (2.2,2.05) {$P = 2(L + w)$};
\end{tikzpicture}

{\small Dua rumus, yang berlaku untuk setiap persegi panjang: gantilah
nilai $L$ dan $w$ untuk memperoleh bilangan.}
\end{omfigure}

\section{Mengganti nilai}

\begin{method}[Menghitung nilai sebuah bentuk]\label{met:g7:literal:substitute}
Untuk menghitung nilai sebuah bentuk pada nilai hurufnya yang diberikan:
\begin{enumerate}
  \item tulislah ulang bentuknya dengan mengembalikan tanda $\times$
        yang tersembunyi: $3a + a^2$ adalah
        $3 \times a + a \times a$;
  \item gantilah setiap huruf dengan nilainya, dengan memberi tanda
        kurung di sekelilingnya bila perlu;
  \item hitunglah, sambil menuruti urutan pengerjaan pada
        \cref{ch:g7:priorities}.
\end{enumerate}
\end{method}

\begin{example}\label{ex:g7:literal:substitute}
Hitunglah nilai $E = 3x + x^2$ untuk $x = 5$:
\[
E = 3 \times 5 + 5 \times 5 = 15 + 25 = 40 .
\]
Hitunglah nilai $F = 2(a + 3b)$ untuk $a = 4$ dan $b = 0.5$:
\[
F = 2 \times (4 + 3 \times 0.5) = 2 \times (4 + 1.5) = 2 \times 5.5
= 11 .
\]
\end{example}

\section{Menyederhanakan bentuk}

\begin{proposition}[Mengumpulkan suku sejenis]\label{prop:g7:literal:reduce}
Suku yang memuat bagian huruf yang sama dapat dikumpulkan dengan
menjumlahkan bagian bilangannya (inilah sifat distributif,
\cref{thm:g7:priorities:distributivity}, yang dibaca terbalik):
\[
5x + 3x = (5 + 3)x = 8x,
\qquad
7a - 2a = 5a .
\]
Suku dengan bagian huruf yang berbeda (seperti $3x$ dan $5y$, atau $x$
dan $x^2$) tidak dapat dikumpulkan.
\end{proposition}

\begin{example}\label{ex:g7:literal:reduce}
Sederhanakan $E = 4x + 7 + 3x + 2$, dengan memilah sukunya:
\[
E = (4x + 3x) + (7 + 2) = 7x + 9 .
\]
Hati-hati: $7x + 9$ itu \emph{bukan} $16x$! Ujilah dengan $x = 1$:
$4 + 7 + 3 + 2 = 16$ dan $7 \times 1 + 9 = 16$ cocok, tetapi $16x$
juga akan bernilai $16$ --- ujilah dengan $x = 2$: bentuk aslinya
memberi $8 + 7 + 6 + 2 = 23$, dan $7 \times 2 + 9 = 23$, sedangkan
$16 \times 2 = 32$. Penyederhanaan $7x + 9$ benar, $16x$ salah.
\end{example}

\begin{method}[Menguji sebuah kesamaan]\label{met:g7:literal:test}
Untuk memeriksa apakah dua bentuk mungkin sama untuk setiap nilai:
\begin{enumerate}
  \item gantilah dengan nilai yang sama pada keduanya, sedikitnya dua
        kali (misalnya $x = 2$ dan $x = 10$; hindari $0$ dan $1$,
        keduanya menyembunyikan terlalu banyak perbedaan);
  \item kalau hasilnya berbeda walau sekali saja, kedua bentuk itu
        \emph{tidak} sama --- satu contoh penyangkal membunuh
        pernyataan umum;
  \item kalau hasilnya cocok, kesamaannya baru \emph{masuk akal}:
        pembuktiannya memerlukan aljabar (menjabarkan, menyederhanakan).
\end{enumerate}
\end{method}

\begin{example}\label{ex:g7:literal:test}
Apakah $2(x + 3)$ sama dengan $2x + 6$ untuk setiap $x$? Menjabarkan
membuktikannya: $2(x+3) = 2x + 2 \times 3 = 2x + 6$. Apakah
$(x + 1)^2$ sama dengan $x^2 + 1$? Ujilah $x = 2$: $(2+1)^2 = 9$
tetapi $2^2 + 1 = 5$. Bukan --- dan satu contoh penyangkal sudah
menjadi bukti lengkap untuk ``bukan''.
\end{example}

\section{Menghasilkan bentuk aljabar}

\begin{example}[Dari keadaan menjadi bentuk]\label{ex:g7:literal:produce}
Sebuah taksi memungut $4$ euro ditambah $2$ euro per kilometer. Untuk
perjalanan sejauh $k$ kilometer harganya
\[
p = 4 + 2k \quad\text{euro}.
\]
Untuk $k = 7$: $p = 4 + 14 = 18$ euro. Bentuk itu \emph{adalah} jawaban umumnya;
menggantikan nilai menghasilkan setiap jawaban khususnya.
\end{example}

\begin{example}[Trik bilangan]\label{ex:g7:literal:trick}
``Pikirkan sebuah bilangan, lipatduakan, tambahkan $10$, bagi dengan
$2$, lalu kurangkan bilanganmu tadi.'' Sebutlah bilangan itu $n$ lalu
ikuti langkahnya:
\[
n \ \to\ 2n \ \to\ 2n + 10 \ \to\ \frac{2n + 10}{2} = n + 5 \ \to\
n + 5 - n = 5 .
\]
Semua orang memperoleh $5$, berapa pun $n$ tadi --- hurufnya
menjelaskan sihirnya.
\end{example}

\section{Latihan}

\begin{exercise}[$\star$]\label{exo:g7:literal:1}
Sederhanakan penulisannya:
\[
7 \times x, \qquad
y \times 3, \qquad
a \times b \times 5, \qquad
x \times x, \qquad
2 \times (a + 4).
\]
\end{exercise}

\begin{exercise}[$\star$]\label{exo:g7:literal:2}
Kembalikan tanda $\times$-nya:
\[
5ab, \qquad
3(x + 2), \qquad
x^2 + 4x, \qquad
2a^3 .
\]
\end{exercise}

\begin{exercise}[$\star$]\label{exo:g7:literal:3}
Hitunglah nilainya untuk $x = 3$: \ $5x + 1$; \quad $x^2$; \quad
$2x^2 - x$; \quad $10 - 2x$.
\end{exercise}

\begin{exercise}[$\star$]\label{exo:g7:literal:4}
Hitunglah nilai $3a + 2b$ dan $a b + 5$ untuk $a = 4$, $b = 1.5$.
\end{exercise}

\begin{exercise}[$\star$]\label{exo:g7:literal:5}
Sederhanakanlah:
\[
8x + 5x, \qquad
9a - 4a + a, \qquad
6y + 2 + 3y + 7, \qquad
5t + 3 - 2t - 3 .
\]
\end{exercise}

\begin{exercise}[$\star$]\label{exo:g7:literal:6}
Berikan rumus untuk: \omterm{def:g6:measure:perimeter}{keliling} persegi dengan sisi $c$; \omterm{def:g6:measure:perimeter}{keliling}
\omterm{def:g6:shapes:triangles}{segitiga sama kaki} dengan alas $b$ dan kaki $s$; harga $n$ buah roti
seharga $1.20$ euro masing-masing.
\end{exercise}

\begin{exercise}[$\star$]\label{exo:g7:literal:7}
Pasangkan setiap bentuk dengan gambarannya: $2n$; $n + 2$; $n^2$;
$\frac n2$. Gambarannya: ``setengah dari $n$''; ``dua kali $n$'';
``kuadrat dari $n$''; ``dua lebih banyak daripada $n$''.
\end{exercise}

\begin{exercise}[$\star\star$]\label{exo:g7:literal:8}
Dengan \cref{met:g7:literal:test}, putuskan (dengan bukti atau contoh
penyangkal) apakah kesamaan berikut berlaku untuk setiap $x$:
\[
3(x + 4) = 3x + 12; \qquad
2x + 5 = 7x; \qquad
x + x = x^2 .
\]
\end{exercise}

\begin{exercise}[$\star\star$]\label{exo:g7:literal:9}
Sebuah persegi panjang berlebar $w$ dan panjangnya $3$ cm lebih
daripada lebarnya.
\begin{enumerate}
  \item Nyatakan panjangnya, lalu \omterm{def:g6:measure:perimeter}{kelilingnya}, dalam $w$, lalu
        sederhanakan.
  \item Hitunglah \omterm{def:g6:measure:perimeter}{kelilingnya} untuk $w = 5$ cm, secara langsung dan
        dengan rumusmu yang sudah disederhanakan. Jawabannya sama?
\end{enumerate}
\end{exercise}

\begin{exercise}[$\star\star$]\label{exo:g7:literal:10}
Ikutilah trik pada \cref{ex:g7:literal:trick} dengan huruf: ``pikirkan
sebuah bilangan, tambahkan $6$, kalikan dengan $3$, kurangi $18$, lalu
bagi dengan bilangan awalmu'' (anggap bilangan itu bukan \omterm{def:g1:counting:zero}{nol}). Apa yang
diperoleh semua orang?
\end{exercise}

\begin{exercise}[$\star\star\star$]\label{exo:g7:literal:11}
Gambar di bawah ini adalah persegi dengan sisi $a$ yang pojoknya
dipotong berupa persegi dengan sisi $b$ ($b$ lebih kecil daripada
$a$). Tulislah dua bentuk yang berbeda untuk luasnya (satu sebagai
\omterm{ex:g1:subtraction:difference}{selisih}, satu dengan memotong bangun L itu menjadi dua persegi
panjang), lalu periksalah bahwa keduanya cocok untuk $a = 5$, $b = 2$.
\end{exercise}

\section{Soal: Peramal batang korek api}

\begin{problem}[{Soal akhir pekan --- pola yang bertumbuh dan suku
ke-$n$-nya: tiga rumus untuk satu pola, dan pola terkenal yang patah
pada langkah keenam}]
\label{pb:g7:literal:1}
Susunlah batang korek api menjadi sebaris persegi: satu persegi, lalu
dua, lalu tiga\,\dots\ Berapa batang yang diperlukan gambar yang
keseratus? Membilangnya satu per satu akan memakan seluruh akhir
pekan; sebuah \emph{rumus dalam $n$} menjawab seketika untuk setiap
gambar sekaligus --- persis itulah gunanya huruf
(\cref{ex:g7:literal:formulas}). Tetapi hati-hati: soal ini ditutup
dengan pola termasyhur yang bermain jujur selama lima langkah lalu
mengkhianati semua orang yang mempercayainya.

\medskip
\noindent\textbf{Bagian I --- Sebaris persegi.}
Gambar $1$ adalah satu persegi dari batang korek api ($4$ batang);
setiap gambar berikutnya menambah satu persegi lagi yang berbagi sisi
dengan persegi sebelumnya.
\begin{enumerate}
  \item Gambarlah gambar $1$ sampai $4$ lalu bilanglah batangnya.
  \item Jelaskan mengapa setiap persegi baru memakan tepat $3$ batang
        tambahan, lalu pakailah hal itu untuk meramalkan hitungan
        gambar $10$ tanpa rumus apa pun.
  \item Berilah alasan untuk rumusnya: gambar $n$ memerlukan $3n + 1$
        batang. (Dari mana datangnya ``$+1$'' yang kesepian itu?)
        Periksalah dengan pertanyaan 1.
  \item Amir bernalar dengan cara lain: ``$4$ batang untuk persegi
        yang pertama, lalu $3$ batang untuk masing-masing $n - 1$
        persegi yang lain'', sehingga memberi $4 + 3(n - 1)$.
        Jabarkan lalu sederhanakan bentuknya
        (\cref{prop:g7:literal:reduce}): apakah ia sepakat dengan
        pertanyaan 3?
  \item Lina membilang dengan cara ketiga: ``$2n$ batang mendatar, dan
        $n + 1$ batang tegak.'' Tulislah rumusnya, sederhanakan, lalu
        simpulkan. Setelah itu jawablah pertanyaan pembukanya: berapa
        batang untuk gambar $100$?
\end{enumerate}

\medskip
\noindent\textbf{Bagian II --- Segitiga, meja, dan berjalan mundur.}
\begin{enumerate}[resume]
  \item Sebaris segitiga dari batang korek api (setiap segitiga baru
        bersandar pada yang sebelumnya): bilanglah batangnya untuk
        $1$, $2$, $3$ segitiga, lalu berilah alasan untuk rumus
        $2n + 1$ pada $n$ segitiga.
  \item Berjalan mundur: dengan tepat $49$ batang, berapa segitiga
        yang dapat dibuat sebaris? (Bongkarlah rumusnya selangkah
        demi selangkah.)
  \item Meja jamuan berbentuk persegi menampung satu tamu pada tiap
        sisi yang kosong. Ketika $n$ meja dirapatkan menjadi satu
        baris panjang, jelaskan mengapa rumus tempat duduknya
        $2n + 2$, lalu periksalah untuk $1$, $2$, $3$ meja.
  \item Sebuah pesta menantikan $30$ tamu. Berapa meja yang
        diperlukan satu baris? Penyedia jamuannya malah membelah meja
        yang sama menjadi \emph{dua} baris berisi tujuh: berapa tamu
        yang muat sekarang? Jelaskan keuntungannya dengan rumusnya.
  \item Pola titik: gambar $n$ adalah huruf L dari titik, setinggi $n$
        titik dan selebar $n$ titik (titik pojoknya dibilang sekali).
        Bilanglah titik gambar $1$, $2$, $3$, carilah rumusnya, lalu
        namailah bilangan yang dihasilkannya. (Kalau ditumpuk
        bersama, huruf L ini menutup sebuah persegi --- ceritanya
        dilanjutkan pada \cref{pb:g8:equations:1}.)
\end{enumerate}

\medskip
\noindent\textbf{Bagian III --- Percaya, uji, buktikan.}
\begin{enumerate}[resume]
  \item Seorang teman sekelas menyatakan bahwa gambar $n$ pada Bagian
        I memerlukan $3(n + 2) - 5$ batang. Selesaikan pernyataannya
        dengan menjabarkan lalu menyederhanakan --- tanpa perlu diuji.
  \item ``Mengkuadratkan membuat bilangan menjadi lebih besar, jadi
        $n^2 \geq n$ untuk setiap bilangan $n$.'' Ujilah pernyataan
        itu untuk $n = 2$, $3$, $10$; lalu carilah bilangan yang
        membuatnya gagal. Apa yang dibuktikan ketiga uji yang
        berhasil tadi, dan apa yang dibuktikan satu kegagalan itu
        (\cref{met:g7:literal:test})?
  \item Tandailah titik pada sebuah lingkaran lalu gambarlah
        \emph{setiap} \omterm{def:g6:lines:objects}{ruas garis} di antaranya, lalu bilanglah daerah
        di dalam lingkarannya. Kerjakan untuk $2$, $3$, $4$ dan $5$
        titik (letakkan titiknya tidak merata, agar tidak ada tiga
        \omterm{def:g6:lines:objects}{ruas garis} yang berpotongan di satu titik dalam). Catatlah
        hitungannya. Rumus apa yang dibisikkan polanya?
  \item Sekarang $6$ titik, digambar besar dan tidak merata, dibilang
        dengan sangat hati-hati (ada banyak daerah kecil). Berapa
        daerah yang kamu temukan --- dan apa yang terjadi pada
        dugaannya?
  \item Pesan moralnya, dalam dua kalimat: apa yang diperlukan untuk
        \emph{membuktikan} dugaan pelipatduaan itu, dan mengapa rumus
        batang korek api pada Bagian I dan II tidak pernah berada
        dalam bahaya semacam itu? Rayakan dengan ramalan yang
        terbukti: berapa batang untuk sebaris $2\,026$ persegi?

\end{enumerate}
\end{problem}
